\theory{Sturm--Liouville theory}{How can a second-order differential equation reveal the natural modes and frequencies of a system?}{Sturm--Liouville theory studies self-adjoint boundary-value problems, their eigenvalues, eigenfunctions, orthogonality, completeness, and oscillation. It turns differential equations into a spectral decomposition problem.}{Vibrating strings and membranes, heat flow, quantum mechanics, Fourier series, acoustics, and numerical boundary-value problems.}{Self-adjoint differential operators, boundary conditions, and weighted orthogonality produce discrete modes and eigenvalue expansions.}

\paragraph{The problem and history}
Jacques Sturm and Joseph Liouville developed the theory in the nineteenth century. A regular problem has the form
\[
\frac{d}{dx}\left(p(x)y'\right)+q(x)y=-\lambda w(x)y
\]
on an interval, together with boundary conditions. The coefficient $w$ is a positive weight. The goal is to find values of $\lambda$ for which a nonzero solution exists and the corresponding eigenfunction \cite{sturm_ref1}.

\end{historylist}
\section*{3. Theory}
\subsection*{Core theory}
\paragraph{Eigenvalues and orthogonality}
For a regular problem, the eigenvalues are real, can be ordered
\[
\lambda_1<\lambda_2<\cdots,\qquad \lambda_n\to\infty,
\]
and each has a unique eigenfunction up to multiplication by a constant. Eigenfunctions for different eigenvalues are orthogonal under
\[
\langle f,g\rangle_w=\int_a^b f(x)g(x)w(x)\,dx.
\]
After normalization they form an orthonormal basis, so suitable functions can be expanded as an eigenfunction series \cite{sturm_ref2}.

\paragraph{A model problem and its interpretation}
The equation $-y''=\lambda y$ on $[0,\pi]$ with $y(0)=y(\pi)=0$ has eigenfunctions $\sin(nx)$ and eigenvalues $n^2$. A vibrating string is a sum of these normal modes. Separation of variables converts the heat, wave, and Laplace equations into Sturm--Liouville problems.

The one-dimensional time-independent Schrodinger equation is a Sturm--Liouville problem. Its eigenvalues are energy levels and eigenfunctions are stationary states. Bessel and Legendre equations can also be written in Sturm--Liouville form, explaining the orthogonality of special functions \cite{sturm_ref3}.

\paragraph{Singular problems and computation}
When an endpoint is infinite or a coefficient vanishes, the problem is singular. Weyl's limit-point and limit-circle alternative determines whether an additional boundary condition is needed. Numerical methods include shooting, finite differences, Rayleigh--Ritz, matrix variational methods, and spectral parameter power series.

\section*{4. Important theorems and results}
\begin{enumerate}[leftmargin=*,itemsep=0.35em]
\item \textbf{Reality and ordering.} Regular Sturm--Liouville eigenvalues are real, simple, and increase without bound.

\item \textbf{Orthogonality.} Eigenfunctions belonging to distinct eigenvalues are orthogonal with the weight $w$.

\item \textbf{Completeness.} Normalized eigenfunctions form an orthonormal basis of the weighted $L^2$ space.

\item \textbf{Oscillation theorem.} The $n$th eigenfunction has a controlled number of zeros in the interval.

\item \textbf{Rayleigh quotient principle.} The smallest eigenvalue minimizes an energy quotient among admissible nonzero functions.

\item \textbf{Weyl alternative.} Singular endpoints are classified as limit-point or limit-circle, determining boundary data.
\end{enumerate}
\noindent\emph{Source for these results:} \cite{sturm_ref1}.

\theoryapplications
\theoryconnections{sturm_liouville_theory}{%
\item An ordinary differential equation with boundary conditions becomes a self-adjoint operator problem.
\item Integration supplies the weighted inner product, Hilbert spaces make orthogonality and convergence precise, and spectral theory proves that the allowed eigenvalues are the natural frequencies of the problem.
\item The boundary conditions are essential: they select which modes fit the physical domain and make the operator's spectrum discrete in the standard compact cases \cite{sturm_ref1,sturm_ref2}.
}{%
\item The eigenfunctions form a basis for expanding a temperature, displacement, or wave profile, so a PDE can be solved one mode at a time.
\item In quantum mechanics they are stationary states; in acoustics they are resonant shapes; in numerical analysis they provide basis functions whose eigenvalues reveal stiffness and stability.
\item The connection works because self-adjointness makes different modes orthogonal and keeps the expansion coefficients independently computable \cite{sturm_ref1,sturm_ref3}.
}
\section*{Glossary}
\begin{description}[leftmargin=*,style=nextline,itemsep=0.3em]
\item[\textbf{Sturm–Liouville problem.}] A second-order differential equation with boundary conditions that leads to a self-adjoint eigenvalue problem.
\item[\textbf{Eigenvalue.}] A value $\lambda$ for which the differential equation has a nonzero solution satisfying the boundary conditions.
\item[\textbf{Eigenfunction.}] A nonzero solution corresponding to an eigenvalue; these form orthogonal basis functions.
\item[\textbf{Weight function.}] A positive function $w(x)$ that appears in the differential equation and defines the inner product for orthogonality.
\item[\textbf{Self-adjoint operator.}] An operator equal to its adjoint; ensures real eigenvalues and orthogonal eigenfunctions.
\item[\textbf{Boundary conditions.}] Constraints at the endpoints that select which solutions are physically admissible.
\item[\textbf{Orthogonality.}] The property that eigenfunctions for different eigenvalues integrate to zero with respect to the weight function.
\item[\textbf{Singular problem.}] A Sturm–Liouville problem where an endpoint is infinite or a coefficient vanishes, requiring special analysis.
\item[\textbf{Rayleigh quotient.}] The energy-like ratio $\langle f,Tf\rangle/\langle f,f\rangle$ achieving its minimum at the smallest eigenvalue.
\item[\textbf{Weyl alternative.}] The classification of singular endpoints as limit-point or limit-circle, determining whether additional boundary conditions are needed.
\item[\textbf{Separation of variables.}] A method of solving PDEs by writing solutions as products of single-variable functions.
\end{description}
\section*{7. Sample Questions}
\emph{Exercise reference:} \cite{sturm_ref1}. The cited edition is the source for the exercise set; consult its Exercises section for the official statement and numbering.
\begin{enumerate}[leftmargin=*,itemsep=0.9em]
\item \textbf{Exercise 1.} Solve the Sturm--Liouville problem $-y'' = \lambda y$ on $[0, \pi]$ with $y(0) = 0$ and $y'(\pi) = 0$. Find the first four eigenvalues. \emph{Solution.} For $\lambda > 0$, write $\lambda = \mu^2$. The general solution is $y = A\sin(\mu x) + B\cos(\mu x)$. From $y(0) = 0$: $B = 0$. From $y'(\pi) = 0$: $A\mu\cos(\mu\pi) = 0$, so $\mu\pi = (n - \frac{1}{2})\pi$, giving $\mu = n - \frac{1}{2}$ for $n = 1, 2, 3, \ldots$ The eigenvalues are $\lambda_n = (n - \frac{1}{2})^2$. The first four are $\frac{1}{4}$, $\frac{9}{4}$, $\frac{25}{4}$, and $\frac{49}{4}$.

\item \textbf{Exercise 2.} Verify that the eigenfunctions $y_n(x) = \sin\!\left((n - \frac{1}{2})x\right)$ and $y_m(x) = \sin\!\left((m - \frac{1}{2})x\right)$ with $n \neq m$ are orthogonal on $[0, \pi]$ with weight $w(x) = 1$. \emph{Solution.} Compute $\int_0^\pi \sin\!\left((n - \frac{1}{2})x\right)\sin\!\left((m - \frac{1}{2})x\right)dx$. Using the product-to-sum formula: $\frac{1}{2}\int_0^\pi \left[\cos\!\left((n - m)x\right) - \cos\!\left((n + m - 1)x\right)\right]dx$. For $n \neq m$ (both integers), $(n - m)$ and $(n + m - 1)$ are both nonzero integers. Then $\int_0^\pi \cos(kx)\,dx = \frac{\sin(k\pi)}{k} = 0$ for any nonzero integer $k$. Thus the integral is $0$, confirming orthogonality.

\item \textbf{Exercise 3.} For the Sturm--Liouville problem $-y'' = \lambda y$ on $[0,1]$ with $y(0) = y(1) = 0$, compute the Rayleigh quotient $R[f] = \frac{\int_0^1 (f')^2\,dx}{\int_0^1 f^2\,dx}$ for $f(x) = x(1-x)$, and compare it with the smallest eigenvalue $\lambda_1 = \pi^2$. \emph{Solution.} Compute $\int_0^1 f^2\,dx = \int_0^1 x^2(1-x)^2\,dx = \int_0^1 (x^2 - 2x^3 + x^4)\,dx = \frac{1}{3} - \frac{1}{2} + \frac{1}{5} = \frac{1}{30}$. Compute $f'(x) = 1 - 2x$, so $\int_0^1 (f')^2\,dx = \int_0^1 (1-2x)^2\,dx = \int_0^1 (1 - 4x + 4x^2)\,dx = 1 - 2 + \frac{4}{3} = \frac{1}{3}$. Then $R[f] = \frac{1/3}{1/30} = 10$. Since $\pi^2 \approx 9.87$, we have $R[f] = 10 > \pi^2 = \lambda_1$, consistent with the Rayleigh quotient principle (the minimum is $\lambda_1$, achieved by the first eigenfunction).

\item \textbf{Exercise 4.} Find the eigenvalues and normalized eigenfunctions of $-y'' = \lambda y$ on $[0, L]$ with $y(0) = y(L) = 0$. \emph{Solution.} The general solution for $\lambda = \mu^2 > 0$ is $y = A\sin(\mu x) + B\cos(\mu x)$. From $y(0) = 0$: $B = 0$. From $y(L) = 0$: $\sin(\mu L) = 0$, so $\mu L = n\pi$, giving $\mu_n = \frac{n\pi}{L}$ and $\lambda_n = \frac{n^2\pi^2}{L^2}$ for $n = 1, 2, 3, \ldots$ The normalized eigenfunctions (with $\|y_n\|^2 = 1$) are $\hat{y}_n(x) = \sqrt{\frac{2}{L}}\sin\!\left(\frac{n\pi x}{L}\right)$.

\item \textbf{Exercise 5.} Show that the oscillation theorem holds for the Sturm--Liouville problem $-y'' = \lambda y$ on $[0, \pi]$ with $y(0) = y(\pi) = 0$: verify that the $n$th eigenfunction $\sin(nx)$ has exactly $n - 1$ zeros in the open interval $(0, \pi)$. \emph{Solution.} The $n$th eigenfunction is $y_n(x) = \sin(nx)$. The zeros of $\sin(nx)$ in $(0, \pi)$ are at $x = \frac{k\pi}{n}$ for $k = 1, 2, \ldots, n - 1$. There are exactly $n - 1$ zeros in the open interval. For example, $y_1(x) = \sin(x)$ has $0$ zeros in $(0,\pi)$, $y_2(x) = \sin(2x)$ has $1$ zero at $x = \pi/2$, and $y_3(x) = \sin(3x)$ has $2$ zeros at $x = \pi/3$ and $x = 2\pi/3$. This confirms the oscillation theorem.

\end{enumerate}
\section*{8. References}
\begin{thebibliography}{9}
\bibitem{sturm_ref1} A. Zettl, \emph{Sturm--Liouville Theory}, Cambridge University Press, 2012.
\bibitem{sturm_ref2} R. Courant and D. Hilbert, \emph{Methods of Mathematical Physics}, Wiley-Interscience, 1989.
\bibitem{sturm_ref3} G. Teschl, \emph{Mathematical Methods in Quantum Mechanics}, American Mathematical Society, 2009.
\end{thebibliography}
